\section{Related Work}
\label{sec:related-work}

Prior DeFi security research has established the basic threat landscape but often stops short of an artifact-level bridge between incident archives and replay validation. Early work on flash-loan-driven exploitation demonstrated how composability, low-friction borrowing, and oracle dependence can transform latent logic mistakes into capital-efficient attacks \citep{qin2021attacking}. Systems work on decentralized exchanges and transaction ordering showed that adversarial execution in DeFi must be understood at the level of transaction sequencing and composability rather than isolated contracts \citep{daian2020flashboys}. Broader overviews of DeFi and its security assumptions described the ecosystem as a layered stack of protocols, incentives, and integration risks, highlighting the methodological difficulty of attributing losses to a single locus of failure \citep{werner2021sok,zhou2023sok}.

A second body of work comes from protocol audits and vendor-style security reviews. These analyses are useful because they identify concrete code-level weaknesses, such as access-control gaps, arithmetic precision loss, or unsafe external-call structure. However, audit artifacts also have a known limitation: they mix true exploit preconditions with hardening advice and context-dependent design concerns. A finding that is valid at the code level may never become a replayable exploit under realistic liquidity, oracle, or governance assumptions. Our own repository reflects this tension. The scanner produced 651 findings across five protocols, but manual triage reduced the set to 87 unique issues, illustrating how a research paper built directly on raw scanner output would exaggerate exploitability.

A third line of evidence comes from public incident databases and post-mortems. Resources such as the SlowMist hacked archive and protocol-authored incident analyses are indispensable for historical coverage and case reconstruction \citep{slowmistHacked,eulerPostmortem,bonqIncident,lendfmeIncident}. Yet those sources are heterogeneous. They vary in how they define the target, classify the root cause, estimate loss, and separate direct protocol exploits from ancillary compromise. Our approach uses those archives as the first stage of research, not the last. The paper therefore occupies a middle ground: broader than a single incident post-mortem, narrower than a generic ecosystem survey, and more verification-oriented than a standalone audit report.
